\section*{Bab \ref{ch:b1:diffeq} --- Persamaan Diferensial Linear}

\begin{solution}{exo:b1:diffeq:1}
Homogennya: $y_h = \lambda\,\eu^{-2x}$. Khususnya: cobalah $y_p =
c\,\eu^{3x}$ (karena $3$ bukan akar $r + 2$): $3c + 2c = 1$, $c =
\frac15$. Penyelesaian umumnya: $y = \frac{\eu^{3x}}{5} +
\lambda\,\eu^{-2x}$. Dengan $y(0) = 1$: $\frac15 + \lambda = 1$,
$\lambda = \frac45$, jadi $y = \frac{\eu^{3x} + 4\,\eu^{-2x}}{5}$.
\end{solution}

\begin{solution}{exo:b1:diffeq:2}
Pada $\intoo{0}{+\infty}$, bagilah dengan $x$: $y' - \frac{2}{x}\,y = x^2$.
Di sini $A(x) = -2\ln x$, $\eu^{-A(x)} = x^2$: jadi penyelesaian homogennya
$\lambda x^2$. \omterm{thm:b1:diffeq:voc}{Variasi konstanta}: $\mu'(x) = x^2 \cdot x^{-2} =
1$, jadi $\mu = x + \lambda$ dan
\[
y(x) = x^3 + \lambda x^2, \qquad \lambda \in \R .
\]
\emph{Periksa:} $x(3x^2 + 2\lambda x) - 2(x^3 + \lambda x^2) = x^3$.
\end{solution}

\begin{solution}{exo:b1:diffeq:3}
$y'' - 3y' + 2y = 0$: akarnya $1$ dan $2$; $\;y = \lambda\,\eu^{x} +
\mu\,\eu^{2x}$.

$y'' + 4y' + 4y = 0$: akar gandanya $-2$; $\;y = (\lambda + \mu
x)\,\eu^{-2x}$.

$y'' - 2y' + 5y = 0$: akarnya $1 \pm 2\iu$; $\;y = \eu^{x}(\lambda\cos
2x + \mu\sin 2x)$.
\end{solution}

\begin{solution}{exo:b1:diffeq:4}
Homogennya: akar $\pm 1$, jadi $y_h = \lambda\,\eu^x + \mu\,\eu^{-x}$.
Penyelesaian khusus dengan ruas kanan polinomial ($\gamma = 0$ bukan akarnya):
$y_p = ax^2 + bx + c$; dengan mensubstitusikannya, $2a - (ax^2 + bx + c) = x^2$
memberikan $a = -1$, $b = 0$, $c = 2a = -2$: jadi $y_p = -x^2 - 2$. Penyelesaian
umumnya $y = -x^2 - 2 + \lambda\eu^x + \mu\eu^{-x}$.

Cauchy: $y(0) = -2 + \lambda + \mu = 0$ dan $y'(0) = \lambda - \mu =
1$: jadi $\lambda = \frac32$, $\mu = \frac12$. Sehingga
$y = -x^2 - 2 + \frac{3\eu^x + \eu^{-x}}{2}$.
\end{solution}

\begin{solution}{exo:b1:diffeq:5}
$a(x) = \tan x$, $A(x) = -\ln(\cos x)$ (yang sah karena $\cos > 0$ pada
selang itu), $\eu^{-A} = \cos x$: jadi penyelesaian homogennya $\lambda\cos x$.
\omterm{thm:b1:diffeq:voc}{Variasi konstanta}: $\mu'(x) = \sin 2x \cdot \frac{1}{\cos x} =
2\sin x$, jadi $\mu = -2\cos x + \lambda$ dan
\[
y(x) = -2\cos^2 x + \lambda \cos x .
\]
\emph{Periksa:} $y' = 4\cos x \sin x - \lambda\sin x$ dan $y\tan x =
-2\cos x\sin x + \lambda \sin x$; jumlahnya adalah $2\cos x\sin x =
\sin 2x$, sesuai tuntutannya.
\end{solution}

\begin{solution}{exo:b1:diffeq:6}
$\chi(r) = r^2 - 4r + 3 = (r-1)(r-3)$: di sini $\gamma = 1$ akar sederhana
($m = 1$). Cobalah $y_p = x(ax + b)\,\eu^x$. Dengan $u = ax^2 + bx$,
\[
y_p'' - 4y_p' + 3y_p = \bigl(u'' + (2 - 4)u' + \chi(1) u\bigr)\eu^x
= \bigl(2a - 2(2ax + b)\bigr)\eu^x .
\]
Samakan dengan $(2x + 1)\eu^x$: $-4a = 2$ dan $2a - 2b = 1$, jadi $a =
-\frac12$, $b = -1$. Penyelesaian umumnya:
\[
y = -\Bigl(\frac{x^2}{2} + x\Bigr)\eu^{x} + \lambda\,\eu^{x} +
\mu\,\eu^{3x}, \qquad (\lambda, \mu) \in \R^2 .
\]
\end{solution}

\begin{solution}{exo:b1:diffeq:7}
Homogennya: $y_h = \lambda\cos 2x + \mu\sin 2x$.

Ruas kanan $x$ (dengan $\gamma = 0$ bukan akarnya): $y_1 = ax + b$ dengan
$4(ax + b) = x$: jadi $y_1 = \frac x4$.

Ruas kanan $\sin 2x = \Im(\eu^{2\iu x})$, dengan $2\iu$ akar sederhana
$r^2 + 4$: cobalah $z = c\,x\,\eu^{2\iu x}$; maka $z'' + 4z =
4\iu c\,\eu^{2\iu x}$, yang sama dengan $\eu^{2\iu x}$ untuk $c =
\frac{1}{4\iu} = -\frac{\iu}{4}$. Jadi $z = -\frac{\iu x}{4}(\cos 2x +
\iu \sin 2x)$ dan $y_2 = \Im(z) = -\frac{x\cos 2x}{4}$.

Menurut superposisi:
\[
y = \frac{x}{4} - \frac{x\cos 2x}{4} + \lambda\cos 2x + \mu\sin 2x .
\]
\end{solution}

\begin{solution}{exo:b1:diffeq:8}
Persamaan $T' + kT = 20k$ mempunyai penyelesaian khusus yang konstan $20$ dan
penyelesaian homogen $\lambda\eu^{-kt}$: jadi $T(t) = 20 +
\lambda\,\eu^{-kt}$, dan $T(0) = 80$ memberikan $\lambda = 60$:
\[
T(t) = 20 + 60\,\eu^{-kt} .
\]
$T(10) = 50$: $\eu^{-10k} = \frac12$, jadi $k = \frac{\ln 2}{10}$. Lalu
$T(t) = 25$ menuntut $\eu^{-kt} = \frac{5}{60} = \frac{1}{12}$, yakni
\[
t = \frac{\ln 12}{k} = 10\,\frac{\ln 12}{\ln 2} \approx 35.8
\text{ menit.}
\]
\end{solution}

\begin{solution}{exo:b1:diffeq:9}
\begin{enumerate}
  \item $\chi(r) = r^2 + 2\varepsilon r + 1$, $\Delta = 4(\varepsilon^2
        - 1)$.
        Untuk $\varepsilon \in \intco{0}{1}$: akarnya $-\varepsilon \pm
        \iu\sqrt{1 - \varepsilon^2}$, jadi
        $y = \eu^{-\varepsilon t}\bigl(\lambda\cos\omega t +
        \mu\sin\omega t\bigr)$ dengan $\omega = \sqrt{1 -
        \varepsilon^2}$.
        Untuk $\varepsilon = 1$: akar gandanya $-1$, jadi $y = (\lambda + \mu
        t)\,\eu^{-t}$.
        Untuk $\varepsilon > 1$: akar realnya $r_\pm = -\varepsilon \pm
        \sqrt{\varepsilon^2 - 1}$, keduanya $< 0$, dan $y =
        \lambda\eu^{r_+t} + \mu\eu^{r_-t}$.
  \item Untuk $\varepsilon \in \intoo{0}{1}$: $\abs y \leq
        \eu^{-\varepsilon t}(\abs\lambda + \abs\mu) \to 0$. Untuk
        $\varepsilon = 1$: $(\lambda + \mu t)\eu^{-t} \to 0$
        (karena eksponensial mengalahkan polinomial,
        \cref{prop:b1:functions:powerrules}). Untuk $\varepsilon > 1$:
        kedua eksponensialnya meluruh karena $r_\pm < 0$ (memang
        $\sqrt{\varepsilon^2 - 1} < \varepsilon$). Untuk $\varepsilon =
        0$: $y = \lambda\cos t + \mu\sin t$ mempunyai amplitudo tetap
        $\sqrt{\lambda^2 + \mu^2} \neq 0$ kecuali bila $y = 0$.
  \item Tulis $\lambda\cos\omega t + \mu\sin\omega t =
        R\cos(\omega t - \varphi)$ dengan $R = \sqrt{\lambda^2 + \mu^2}
        > 0$. Nol $y$ adalah nol $\cos(\omega t -
        \varphi)$ (karena faktor $\eu^{-\varepsilon t}$ tak pernah
        lenyap): $\omega t - \varphi \equiv \frac\pi2 \pmod \pi$,
        yaitu barisan aritmetika dengan celah $\frac{\pi}{\omega} =
        \frac{\pi}{\sqrt{1 - \varepsilon^2}}$.
\end{enumerate}
\end{solution}

\begin{solution}{exo:b1:diffeq:10}
Ambil $x = y = 0$: $2f(0) = 2f(0)^2$, dan $f(0) \neq 0$ memberikan $f(0) =
1$. Tetapkan $x$ lalu turunkan persamaannya dua kali terhadap $y$:
\[
f''(x+y) + f''(x-y) = 2 f(x) f''(y) .
\]
Dengan mengambil $y = 0$: $\;2f''(x) = 2 f(x) f''(0)$, yaitu
\[
f''(x) = c\,f(x), \qquad c = f''(0).
\]
\emph{Kasus $c = \omega^2 > 0$:} $f(x) = \lambda\cosh\omega x +
\mu\sinh\omega x$; lalu $f(0) = 1$ memberikan $\lambda = 1$. Dengan memasukkannya ke
persamaan fungsionalnya dan memakai rumus penjumlahan
(\cref{prop:b1:functions:hyprules}), persamaan itu memaksa $\mu = 0$
(bandingkan koefisien $\sinh\omega x \sinh\omega y$ atau
hitung nilainya di $x = y$): jadi $f = \cosh\omega x$, yang memang memenuhi
$\cosh(x+y) + \cosh(x-y) = 2\cosh x\cosh y$.

\emph{Kasus $c = -\omega^2 < 0$:} serupa dengan itu, $f(x) = \cos\omega x$
($\omega \neq 0$), yang memenuhi persamaannya.

\emph{Kasus $c = 0$:} $f$ afin dengan $f(0) = 1$: $f(x) = 1 + \mu x$;
lalu persamaannya memaksa $\mu = 0$, yang tersingkir (karena $f$ tidak konstan).

Kesimpulannya: penyelesaiannya adalah $f(x) = \cos\omega x$ dan $f(x) =
\cosh\omega x$, dengan $\omega > 0$.
\end{solution}

\begin{solution}{exo:b1:diffeq:11}
Tulis $z(t) = y(\eu^t)$, sehingga $y(x) = z(\ln x)$ untuk $x > 0$. Maka
\[
y'(x) = \frac{z'(\ln x)}x,
\qquad
y''(x) = \frac{z''(\ln x) - z'(\ln x)}{x^2} ,
\]
lalu dengan mensubstitusikannya ke dalam persamaannya:
\[
x^2y'' - xy' + y = \bigl(z'' - z'\bigr) - z' + z
= z'' - 2z' + z = 0 .
\]
\omterm{def:b1:diffeq:linear2}{Polinomial karakteristiknya} $(r - 1)^2$: akar gandanya $1$, jadi $z(t) =
(\lambda + \mu t)\,\eu^t$ dan, kembali dalam variabel $x = \eu^t$:
\[
y(x) = (\lambda + \mu\ln x)\,x,
\qquad \lambda, \mu \in \R .
\]
\end{solution}

\begin{solution}{exo:b1:diffeq:12}
\begin{enumerate}
  \item Dalam bentuk ternormalkan $y' - \frac2x\,y = 0$ pada masing-masing selang:
        $A(x) = -2\ln\abs x$, jadi penyelesaiannya adalah $y = a x^2$ pada
        $\intoo0{+\infty}$ dan $y = b x^2$ pada $\intoo{-\infty}0$,
        dengan konstanta yang saling bebas
        (\cref{thm:b1:diffeq:homogeneous1}).
  \item Misalkan $y = ax^2$ untuk $x \geq 0$ dan $bx^2$ untuk $x < 0$. Pada
        masing-masing setengah garis terbukanya, $y$ dapat diturunkan dengan $xy' = 2y$.
        Di $0$: hasil bagi selisihnya $\frac{y(h) - y(0)}h =
        ah$ atau $bh$ menuju $0$, jadi $y'(0) = 0$ ada, dan
        persamaannya di $x = 0$ berbunyi $0 \cdot y'(0) = 2y(0) = 0$:
        yang terpenuhi. Jadi $y$ memenuhi persamaannya pada seluruh $\R$.
  \item Penyelesaian pada $\R$ tepat berupa fungsi rekatan itu: yaitu
        keluarga berparameter \emph{dua} bagi persamaan orde satu.
        Tak ada pertentangan dengan \cref{thm:b1:diffeq:voc},
        yang hipotesisnya gugur di sini: bila ditulis $y' + a(x)y = 0$,
        koefisien $a(x) = -\frac2x$ tidak kontinu di $0$
        --- bahkan tidak terdefinisi --- sehingga $\R$ bukanlah selang
        yang teoremanya berlaku di sana. Kesingularan di $0$
        memutuskan kedua setengah garisnya, dan nilai $y(0) = 0$
        terpaksa demikian tanpa membawa informasi menyeberang. Setiap data
        Cauchy di $x_0 \neq 0$ hanya menentukan penyelesaiannya pada
        setengah garis yang memuat $x_0$.
\end{enumerate}
\end{solution}

\begin{solution}{pb:b1:diffeq:1}
\textbf{1.} $\chi(r) = r^2 + 2\lambda r + \omega_0^2$, $\Delta =
4(\lambda^2 - \omega_0^2) < 0$ untuk $0 < \lambda < \omega_0$: akarnya
$-\lambda \pm \iu\omega_d$, jadi menurut
\cref{thm:b1:diffeq:homogeneous2}
\[
x(t) = \eu^{-\lambda t}\bigl(\lambda_1\cos\omega_d t +
\mu_1\sin\omega_d t\bigr), \qquad (\lambda_1, \mu_1) \in \R^2 .
\]
Untuk $\lambda = 0$: $x = \lambda_1\cos\omega_0 t +
\mu_1\sin\omega_0 t$. Regim kritis ($\lambda = \omega_0$) dan
teredam lebih ($\lambda > \omega_0$) adalah regim pada
\cref{exo:b1:diffeq:9} (setelah waktunya diskalakan ulang): yaitu $(\lambda_1 +
\mu_1 t)\eu^{-\lambda t}$, dan kombinasi $\eu^{r_\pm t}$
dengan $r_\pm = -\lambda \pm \sqrt{\lambda^2 - \omega_0^2}$.

\textbf{2.} Teredam lemah: $\abs x \leq \eu^{-\lambda
t}(\abs{\lambda_1} + \abs{\mu_1}) \to 0$. Kritis: $(\lambda_1 +
\mu_1 t)\eu^{-\lambda t} \to 0$ karena eksponensial mengalahkan
polinomial (\cref{prop:b1:functions:powerrules}). Teredam lebih:
$r_- < r_+ = -\lambda + \sqrt{\lambda^2 - \omega_0^2} < 0$
karena $\sqrt{\lambda^2 - \omega_0^2} < \lambda$; jadi kedua
eksponensialnya meluruh.

\textbf{3.} Sepanjang sebuah penyelesaian $(H)$, dengan memakai $x'' = -2\lambda x' -
\omega_0^2 x$:
\[
\mathcal E'(t) = x'x'' + \omega_0^2 x x'
= x'\bigl(-2\lambda x' - \omega_0^2 x\bigr) + \omega_0^2 xx'
= -2\lambda\,x'^2 \leq 0 .
\]
Jika $x(t_0) = x'(t_0) = 0$ maka $\mathcal E(t_0) = 0$; di sini $\mathcal E$
tak negatif dan tak naik, sehingga $\mathcal E \equiv 0$ pada
$\intco{t_0}{+\infty}$, yang memaksa $x \equiv 0$ di sana; untuk $t \leq
t_0$, jalankan argumen yang sama pada $\tilde x(t) = x(2t_0 - t)$, yang
memenuhi persamaan dengan redaman $-\lambda$ tetapi tetap mempunyai
$\tilde{\mathcal E}(t_0) = 0$ dan $\tilde{\mathcal E}' =
+2\lambda\tilde x'^2 \geq 0$ dengan $\tilde{\mathcal E} \geq 0$;
tak negatif, tidak turun, dan nol di ujung kanan
$\intoc{-\infty}{t_0}$ berarti nol di sepanjangnya. Jadi $x \equiv 0$ pada
$\R$ --- yaitu bukti ketunggalan lewat energi, yang sah untuk setiap $\lambda
\geq 0$.

\textbf{4.} $x(t + T_d) = R\,\eu^{-\lambda t}\eu^{-\lambda
T_d}\cos(\omega_d t + 2\pi - \varphi) = \eu^{-\lambda T_d}x(t)$.
Faktor penyusutan per pseudoperiodenya adalah $\eu^{-\delta}$ dengan
$\delta = \lambda T_d = \frac{2\pi\lambda}{\omega_d}$. Untuk
$\omega_0 = 1$, $\lambda = 0.1$: $\omega_d = \sqrt{0.99} =
0.99499$, jadi $\delta = \frac{0.62832}{0.99499} = 0.6315$: setiap
ayunan menyisakan $\eu^{-0.63} \approx 53\%$ amplitudonya.

\textbf{5.} Faktor amplitudonya adalah $\eu^{-\lambda t}$, yang
meluruh sebesar $\eu$ selama $t = \frac1\lambda$. Selang itu memuat
$\frac{1/\lambda}{T_d} = \frac{\omega_d}{2\pi\lambda}$
pseudoperiode. Untuk $\lambda \ll \omega_0$, $\omega_d \approx
\omega_0$ sehingga ini $\approx \frac{\omega_0}{2\pi\lambda} =
\frac Q\pi$. Senar gitar ber-$Q = 300$ berdenting sekitar seratus
periode; sedangkan peredam pintu dengan $Q = 1$ tak menuntaskan satu pun.

\textbf{6.} Mensubstitusikan $z\,\eu^{\iu\Omega t}$ ke ruas
kirinya memberikan $z\,(-\Omega^2 + 2\iu\lambda\Omega +
\omega_0^2)\,\eu^{\iu\Omega t}$, yang sama dengan $A\,\eu^{\iu\Omega
t}$ tepat untuk $z = \frac{A}{\omega_0^2 - \Omega^2 +
2\iu\lambda\Omega}$ (penyebutnya tak nol: bagian imajinernya
adalah $2\lambda\Omega > 0$). Karena koefisiennya real,
bagian real $x_p = \Re\bigl(z\eu^{\iu\Omega t}\bigr)$ memenuhi
persamaan dengan ruas kanan $\Re\bigl(A\eu^{\iu\Omega t}\bigr) =
A\cos\Omega t$.

\textbf{7.} Tulis $\omega_0^2 - \Omega^2 + 2\iu\lambda\Omega =
\sqrt{D}\,\eu^{\iu\varphi}$ dengan $D = (\omega_0^2 - \Omega^2)^2 +
4\lambda^2\Omega^2$ dan $\varphi \in \intoo0\pi$ (karena bagian
imajinernya $2\lambda\Omega$ positif), sehingga $\tan\varphi =
\frac{2\lambda\Omega}{\omega_0^2 - \Omega^2}$. Lalu $z =
\frac{A}{\sqrt D}\eu^{-\iu\varphi}$ dan
\[
x_p(t) = \Re\Bigl(\frac A{\sqrt D}\,\eu^{\iu(\Omega t -
\varphi)}\Bigr) = R\cos(\Omega t - \varphi),
\qquad R = \frac A{\sqrt D} .
\]

\textbf{8.} Ketika $\Omega \to 0^+$: $D \to \omega_0^4$, jadi $R \to
A/\omega_0^2$ dan $\tan\varphi \to 0^+$ dengan $\varphi \in
\intoo0{\frac\pi2}$: sehingga $\varphi \to 0$. Massanya mengikuti gayanya
secara kuasi-statis, tergeser sebesar gaya dibagi kekakuan. Ketika $\Omega \to
+\infty$: $D \sim \Omega^4$, jadi $R \sim A/\Omega^2 \to 0$, dan
$\varphi \to \pi$ (bilangan kompleks $\omega_0^2 - \Omega^2 +
2\iu\lambda\Omega$ menuju kuadran kedua dengan argumen $\to
\pi$): jadi massanya nyaris tak bergerak, dan berlawanan fase ---
karena kelembamannya mendominasi.

\textbf{9.} Menurut \cref{thm:b1:diffeq:structure2}~(1), setiap penyelesaian
berupa $x = x_p + x_h$ dengan $x_h$ memenuhi $(H)$; menurut pertanyaan 2, $x_h(t)
\to 0$, jadi $x(t) - x_p(t) \to 0$: semua penyelesaian konvergen ke
keadaan tunak yang sama. Syarat awalnya hanya membentuk peralihannya.

\textbf{10.} Jika $x$ penyelesaian periodik, maka $x - x_p = x_h$ adalah
penyelesaian periodik $(H)$ yang menuju $0$ di $+\infty$; dan
fungsi periodik yang berlimit $0$ identik dengan $0$ (karena nilainya pada
satu periode berulang selamanya, sehingga setiap nilainya menjadi limit sebuah
barisan bagian yang menuju $0$). Jadi $x = x_p$.

\textbf{11.} Di sini $\lambda = 1$, $\omega_0^2 = 2$, $\Omega = 1$,
$A = 1$: jadi $z = \frac1{2 - 1 + 2\iu} = \frac{1 - 2\iu}5$, sehingga
\[
x_p = \Re\Bigl(\frac{(1 - 2\iu)(\cos t + \iu\sin t)}5\Bigr)
= \frac{\cos t + 2\sin t}5 .
\]
Homogennya: akar $r^2 + 2r + 2$ adalah $-1 \pm \iu$: jadi $x_h =
\eu^{-t}(C\cos t + S\sin t)$. Syaratnya: $x(0) = \frac15 + C = 0$
memberikan $C = -\frac15$; lalu dengan menurunkannya, $x'(0) = \frac25 - C + S =
0$ memberikan $S = C - \frac25 = -\frac35$. Jadi
\[
x(t) = \underbrace{\frac{\cos t + 2\sin t}5}_{\text{tunak}}
- \underbrace{\eu^{-t}\,\frac{\cos t + 3\sin
t}5}_{\text{peralihan}} ,
\]
dengan peralihannya mati seperti $\eu^{-t}$.

\textbf{12.} Dengan menjabarkannya, $g(u) = u^2 - 2(\omega_0^2 - 2\lambda^2)u
+ \omega_0^4 = (u - u_r)^2 + g(u_r)$ dengan $u_r = \omega_0^2 -
2\lambda^2$ dan
\[
g(u_r) = \omega_0^4 - u_r^2 = (\omega_0^2 - u_r)(\omega_0^2 +
u_r) = 2\lambda^2\bigl(2\omega_0^2 - 2\lambda^2\bigr) =
4\lambda^2(\omega_0^2 - \lambda^2) .
\]
Jika $2\lambda^2 < \omega_0^2$, maka $u_r > 0$ merupakan kuadrat
frekuensi yang sah: di sana $g$ mempunyai minimum tegas, sehingga $R = A/\sqrt
g$ mempunyai maksimum tegas di $\Omega_r = \sqrt{u_r} =
\sqrt{\omega_0^2 - 2\lambda^2}$, dengan $R_{\max} = A/\sqrt{g(u_r)}
= \frac{A}{2\lambda\sqrt{\omega_0^2 - \lambda^2}}$.

\textbf{13.} $R(0) = A/\omega_0^2$, jadi
\[
\frac{R_{\max}}{R(0)}
= \frac{\omega_0^2}{2\lambda\sqrt{\omega_0^2 - \lambda^2}}
= \frac{\omega_0}{2\lambda}\cdot
\frac{\omega_0}{\sqrt{\omega_0^2 - \lambda^2}}
= Q\,\Bigl(1 - \frac{\lambda^2}{\omega_0^2}\Bigr)^{-1/2}
\geq Q .
\]
Untuk redaman lemah faktor koreksinya dekat dengan $1$: jadi \omterm{ex:b1:diffeq:oscillation}{resonansi}
mengalikan pergeseran statisnya dengan $Q$ pada dasarnya.

\textbf{14.} $V(\Omega)^2 = \frac{A^2 u}{g(u)}$ dengan $u =
\Omega^2$. Turunannya bertanda sama dengan $g(u) - u\,g'(u) =
(\omega_0^2 - u)^2 + 4\lambda^2 u - u\bigl(2(u - \omega_0^2) +
4\lambda^2\bigr) = (\omega_0^2 - u)^2 + 2u(\omega_0^2 - u) =
(\omega_0^2 - u)(\omega_0^2 + u)$, yang positif untuk $u < \omega_0^2$
dan negatif setelahnya: jadi maksimum tegasnya tepat di $\Omega =
\omega_0$, untuk redaman apa pun. Di sana $\tan\varphi$ meledak dengan
$\varphi \in \intoo0\pi$: jadi $\varphi = \frac\pi2$, dan $x_p'(t) =
-R\Omega\sin(\Omega t - \frac\pi2) = R\Omega\cos(\Omega t)$ tepat
sefase dengan gayanya: yaitu pemindahan daya yang optimal.

\textbf{15.} $R = R_{\max}/\sqrt2 \iff g(u) = 2g(u_r) \iff (u -
u_r)^2 = g(u_r) = 4\lambda^2(\omega_0^2 - \lambda^2)$, yang memberikan
\[
u_\pm = u_r \pm 2\lambda\sqrt{\omega_0^2 - \lambda^2},
\qquad
\Omega_+^2 - \Omega_-^2 = 4\lambda\sqrt{\omega_0^2 - \lambda^2} .
\]
Lalu $\Omega_+ - \Omega_- = \frac{\Omega_+^2 -
\Omega_-^2}{\Omega_+ + \Omega_-}$, dan untuk $\lambda \ll \omega_0$
kedua $\Omega_\pm \approx \omega_0$: sehingga $\Omega_+ - \Omega_- \approx
\frac{4\lambda\omega_0}{2\omega_0} = 2\lambda$, jadi
$\frac{\omega_0}{\Omega_+ - \Omega_-} \approx
\frac{\omega_0}{2\lambda} = Q$. Jadi mengukur lebar puncak \omterm{ex:b1:diffeq:oscillation}{resonansi}
berarti mengukur faktor kualitasnya.

\textbf{16.} $Q = 10$; $\Omega_r = \sqrt{1 - 2(0.05)^2} =
\sqrt{0.995} = 0.9975$; $R_{\max} = \frac1{2 \times 0.05
\sqrt{1 - 0.0025}} = \frac1{0.1 \times 0.99875} = 10.01$; tanggapan
statisnya $R(0) = 1$; lebar pitanya $\approx 2\lambda = 0.1$. Jadi sebuah
paku yang tinggi dan tipis setinggi $\approx Q$ di atas dataran setinggi $1$.

\textbf{17.} Jika $2\lambda^2 \geq \omega_0^2$ maka $u_r \leq 0$ dan
$g'(u) = 2(u - u_r) > 0$ untuk setiap $u > 0$: jadi $g$ naik tegas
pada $\intoo0{+\infty}$, sehingga $R = A/\sqrt g$ turun tegas dari
$R(0) = A/\omega_0^2$: tanggapannya terbesar pada frekuensi nol
dan tak ada puncak sama sekali.

\textbf{18.} Di sini $\gamma = \iu\Omega$ bukan akar $r^2 +
\omega_0^2$ (karena $\Omega \neq \omega_0$), jadi
\cref{met:b1:diffeq:particular} dengan $m = 0$ memberikan $x_p =
\frac{A\cos\Omega t}{\omega_0^2 - \Omega^2}$ (substitusikan lalu
periksa: $-\Omega^2 + \omega_0^2$ dikali kosinusnya). Penyelesaian
umumnya:
\[
x(t) = \frac{A\cos\Omega t}{\omega_0^2 - \Omega^2}
+ \lambda_1\cos\omega_0 t + \mu_1\sin\omega_0 t .
\]

\textbf{19.} Syarat $x(0) = 0$ memaksa $\lambda_1 = -\frac A{\omega_0^2 -
\Omega^2}$, dan $x'(0) = 0$ memaksa $\mu_1 = 0$:
\[
x(t) = \frac{A}{\omega_0^2 - \Omega^2}\,
\bigl(\cos\Omega t - \cos\omega_0 t\bigr)
= \frac{2A}{\omega_0^2 - \Omega^2}\,
\sin\Bigl(\frac{(\omega_0 - \Omega)t}2\Bigr)
\sin\Bigl(\frac{(\omega_0 + \Omega)t}2\Bigr) ,
\]
menurut rumus hasil kali $\cos a - \cos b =
2\sin\frac{b - a}2\sin\frac{b + a}2$ yang diterapkan dengan $a = \Omega t$,
$b = \omega_0 t$.

\textbf{20.} Untuk $\Omega$ yang dekat dengan $\omega_0$, sinus kedua
berayun pada frekuensi cepat $\frac{\omega_0 + \Omega}2
\approx \omega_0$, sedangkan yang pertama menjadi selubung lambat berfrekuensi
$\frac{\abs{\omega_0 - \Omega}}2$: jadi amplitudo osilasi
cepatnya membesar dan mengecil dengan periode selubung
$\frac{2\pi}{\abs{\omega_0 - \Omega}}$ (dua \omterm{pb:b1:diffeq:1}{layangan} per periode
selubung), dan mencapai maksimum $\frac{2A}{\abs{\omega_0^2 - \Omega^2}}$.
Ketika $\Omega \to \omega_0$, \omterm{pb:b1:diffeq:1}{layangannya} menjadi lebih lambat (periodenya
$\to \infty$) sekaligus lebih tinggi (amplitudonya $\to \infty$).

\textbf{21.} Tetapkan $t$. Ketika $\Omega \to \omega_0$:
\[
\frac{2A}{\omega_0^2 - \Omega^2}\sin\Bigl(\frac{(\omega_0 -
\Omega)t}2\Bigr)
= \frac{2A}{\omega_0 + \Omega}\cdot
\frac{\sin\bigl(\frac{(\omega_0 - \Omega)t}2\bigr)}
{\omega_0 - \Omega}
\longrightarrow \frac{2A}{2\omega_0}\cdot\frac t2
= \frac{At}{2\omega_0} ,
\]
sedangkan $\sin\bigl(\frac{(\omega_0 + \Omega)t}2\bigr) \to
\sin(\omega_0 t)$: jadi limitnya adalah $x_\infty(t) =
\frac{At\sin\omega_0 t}{2\omega_0}$. Pemeriksaan langsung: dengan $C =
\frac A{2\omega_0}$, $x_\infty = Ct\sin\omega_0 t$ mempunyai
$x_\infty'' = 2C\omega_0\cos\omega_0 t - C\omega_0^2 t
\sin\omega_0 t$, sehingga $x_\infty'' + \omega_0^2 x_\infty =
2C\omega_0\cos\omega_0 t = A\cos\omega_0 t$, dengan $x_\infty(0) =
0$ dan $x_\infty'(0) = C\sin 0 + C\omega_0 \cdot 0 \cdot \cos 0 =
0$ --- dan untuk $\omega_0 = A = 1$ inilah
persis \cref{ex:b1:diffeq:oscillation}. Jadi \omterm{ex:b1:diffeq:oscillation}{resonansi} adalah
kemerosotan \omterm{pb:b1:diffeq:1}{layangan}: yaitu gelembung pertama selubungnya, yang direntangkan
sampai panjang tak hingga.

\textbf{22.} Tanpa redaman, amplitudo resonannya tumbuh linear
dan tanpa batas; dengan redaman $\lambda > 0$ pertumbuhannya jenuh
pada $R_{\max} \approx Q\,\frac A{\omega_0^2}$. Mekanismenya:
pelesapan mengambil energi dengan laju yang tumbuh bersama amplitudonya
(pertanyaan 23), sehingga penumpukannya berhenti tepat ketika redamannya
membakar energi secepat pemaksaannya memasoknya.

\textbf{23.} Dengan $x_p' = -R\Omega\sin(\Omega t - \varphi)$, sepanjang
satu periode rata-ratanya $\langle\cos^2\rangle =
\langle\sin^2\rangle = \frac12$ dan $\langle\sin\cos\rangle = 0$
memberikan:
\[
\langle P_{\mathrm{diss}}\rangle
= 2\lambda\,R^2\Omega^2\,\langle\sin^2(\Omega t -
\varphi)\rangle = \lambda R^2\Omega^2 ;
\]
lalu dengan menjabarkan $\sin(\Omega t - \varphi) = \sin\Omega t\cos\varphi
- \cos\Omega t\sin\varphi$:
\[
\langle P_{\mathrm{in}}\rangle
= -AR\Omega\,\bigl\langle\cos\Omega t\,\sin(\Omega t -
\varphi)\bigr\rangle
= AR\Omega\,\frac{\sin\varphi}2 .
\]
Karena $\sin\varphi = \frac{2\lambda\Omega}{\sqrt D} =
\frac{2\lambda\Omega R}A$, ini sama dengan $\frac{AR\Omega}2 \cdot
\frac{2\lambda\Omega R}A = \lambda R^2\Omega^2$: jadi daya yang disuntikkan dan
yang dilesapkan berimbang persis --- yaitu sifat yang mendefinisikan
regim tunak.

\textbf{24.} (i) Teorema strukturnya memilah setiap penyelesaian menjadi
keadaan tunak ditambah peralihan (pertanyaan 9--11) dan menyusutkan
ketunggalannya menjadi masalah homogen. (ii) Metode kompleksnya
mengubah pencarian penyelesaian khusus menjadi satu pembagian
bilangan kompleks (pertanyaan 6), dengan amplitudo dan fasenya terbaca dari
sebuah \omterm{def:b1:complex:field}{modulus} dan sebuah argumen. (iii) Kurva \omterm{ex:b1:diffeq:oscillation}{resonansinya} murni
telaah fungsi --- yaitu sebuah kuadrat dalam $u = \Omega^2$, minimumnya,
\omterm{def:b1:logic:sets}{himpunan} arasnya --- dalam gaya \cref{ch:b1:functions}
(pertanyaan 12--17).

\textbf{25.} Bebas dan teredam: pseudo-osilasi yang meluruh,
dengan umur $\frac1\lambda$ dan sekitar $\frac Q\pi$ ayunan. Terpaksa dan
teredam: peralihannya mati, dan satu keadaan tunak sinusoidal yang tunggal
bertahan pada frekuensi pemaksanya, dengan amplitudo yang memuncak di dekat
$\omega_0$ (tingginya $\approx Q \times$ nilai statisnya, lebarnya $\approx
\frac{\omega_0}Q$) dan fase yang menyapu dari $0$ ke $\pi$ lewat
$\frac\pi2$ di $\omega_0$. Bebas dan tanpa redaman: osilasi
abadi. Terpaksa dan tanpa redaman: \omterm{pb:b1:diffeq:1}{layangan}, yang merosot menjadi
\omterm{ex:b1:diffeq:oscillation}{resonansi} yang tumbuh linear pada penalaan yang persis. Satu bilangan
tanpa dimensi, $Q = \frac{\omega_0}{2\lambda}$, menyetel segalanya ---
tinggi puncak, lebar pita, dan umur peralihannya adalah tiga pembacaan
tombol yang sama. Kelanjutannya: menulis ulang $x'' + 2\lambda x' +
\omega_0^2x$ sebagai sistem orde satu membuka metode matriks pada
\cref{ch:b1:matrices} dan jilid Tahun ke-2, dan menguraikan
pemaksaan periodik yang sebarang menjadi sinusoid (deret Fourier, jilid
Tahun ke-3) menjadikan analisis frekuensi tunggal pada soal ini sebagai
batu bangunan yang universal: selesaikan untuk setiap frekuensi, lalu
superposisikan.

\end{solution}
